% Part:sets-functions-relations
% Chapter: size-of-sets
% Section: reduction-alt

\documentclass[../../../include/open-logic-section]{subfiles}

\begin{document}

\olfileid{sfr}{siz}{red-alt}

\olsection{تقلیل}

\begin{editorial}
  ایہ حصہ تقلیل نال گنتی وار فہرست نہ بنن جوگ ہون نوں ثابت کردا اے،
  تے \olref[nen-alt]{sec} دے نتیجیاں نال میل کھاندا اے۔ اک متبادل،
  ذرا ودھ تفصیلی روپ جیہڑا \olref[nen]{sec} دے نتیجیاں نال میل
  کھاندا اے، \olref[red]{sec} وچ دتا گیا اے۔
\end{editorial}

اسی قطری بناوٹ والی دلیل نال وکھایا سی کہ~$\Bin^\omega$،
!!{nonenumerable} اے۔ اسی اوسے ورگی قطری دلیل نال وکھایا سی کہ
$\Pow{\Nat}$، !!{nonenumerable} اے۔ پر $\Pow{\Nat}$ دے
!!{nonenumerable} ہون نوں ثابت کرن دا اک ہور طریقہ ایہ اے: وکھاؤ کہ
\emph{جے $\Pow{\Nat}$، !!{enumerable} اے تاں $\Bin^\omega$ وی
!!{enumerable} اے}۔ کیوں جے اسی جاندے آں کہ~$\Bin^\omega$،
!!{nonenumerable} اے، ایس لئی ایہ نتیجہ نکلے گا کہ~$\Pow{\Nat}$ وی
ایویں ای اے۔

اینوں اک مسئلے نوں دوجے ول \emph{تقلیل دینا} آکھیا جاندا اے۔ ایس
صورت وچ، اسی~$\Bin^\omega$ دی گنتی وار فہرست بناؤن دے مسئلے نوں
$\Pow{\Nat}$ دی گنتی وار فہرست بناؤن دے مسئلے ول تقلیل دیندے آں۔
پچھلے مسئلے دا حل---$\Pow{\Nat}$ دی اک گنتی وار فہرست---پہلے مسئلے
دا حل---$\Bin^\omega$ دی اک گنتی وار فہرست---دے دیوے گا۔

سیٹ~$B$ دی گنتی وار فہرست بناؤن دے مسئلے نوں سیٹ~$A$ دی گنتی وار
فہرست بناؤن دے مسئلے ول تقلیل دین لئی، اسی~$A$ دی اک گنتی وار فہرست
نوں~$B$ دی گنتی وار فہرست وچ بدلّن دا طریقہ دیندے آں۔ ایہ کرن دا
سب توں سوکھا طریقہ اک !!a{surjection} $f\colon A \to B$ تعریف کرنا
اے۔ جے $x_1$، $x_2$، \dots{}،~$A$ دی گنتی وار فہرست اے، تاں
$f(x_1)$، $f(x_2)$، \dots{}،~$B$ دی گنتی وار فہرست ہووے گی۔ ساڈی
صورت وچ اسی اک !!a{surjection} $f\colon \Pow{\Nat}
\to \Bin^\omega$ لبھ رہے آں۔

\begin{prob}
وکھاؤ کہ جے کوئی !!{injective} فنکشن $g\colon B \to A$ موجود اے،
تے~$B$، !!{nonenumerable} اے، تاں~$A$ وی ایویں ای اے۔ ایہ وکھا کے
کرو کہ تسیں~$g$ نوں~$A$ دی اک گنتی وار فہرست توں~$B$ دی گنتی وار
فہرست بناؤن لئی کِویں ورت سکدے او۔
\end{prob}

\begin{proof}[{تقلیل نال \olref[nen-alt]{thm:nonenum-pownat} دا ثبوت}]
تقلیل لئی فرض کرو کہ~$\Pow{\Nat}$، !!{enumerable} اے، تے ایس لئی
اوہدی اک گنتی وار فہرست $N_{1}$، $N_{2}$، $N_{3}$، \dots اے۔

فنکشن $f \colon \Pow{\Nat} \to \Bin^\omega$ ایویں تعریف کرو کہ
$f(N)$ اوہ لڑی~$s$ ہووے جیہدے لئی $s(n) = 1$ عین اوس ویلے جدوں
$n \in N$، تے ہور صورت وچ $s(n) = 0$ ہووے۔
% OLSIZ-010 ماخذ دی درستی: نتیجے والی لڑی دا ناں s اے؛ انگریزی عبارت وچ غیر متعین k والے s_k تے s_k(n) ورتے گئے سن۔

ایہ صاف طور تے اک فنکشن تعریف کردا اے، کیوں جے جدوں وی
$N \subseteq \Nat$ ہووے، ہر $n \in \Nat$ یا تاں~$N$ دا
!!a{element} اے یا نہیں۔ مثال لئی، جفت قدرتی عدداں دا سیٹ
$2\Nat = \Setabs{2n}{n \in \Nat} = \{0,2, 4, 6,
\dots\}$ لڑی $1010101\dots$ ول، $\emptyset$ لڑی $0000\dots$ ول،
تے~$\Nat$ لڑی $1111\dots$ ول بھیجیا جاندا اے۔

ایہ !!{surjective} وی اے: $0$ تے $1$ دی ہر لڑی قدرتی عدداں دے کسے
سیٹ نال میل کھاندی اے، یعنی اوس سیٹ نال جیہدے رکن اوہ قدرتی عدد نیں
جیہڑے لڑی وچ~$1$ والے مقاماں دے مطابق نیں۔ ہور ٹھیک طور تے، جے
$s \in \Bin^\omega$، تاں $N \subseteq \Nat$ نوں ایویں تعریف کرو:
\[
N = \Setabs{n \in \Nat}{s(n) = 1}
\]
فیر $f(N) = s$، جیویں~$f$ دی تعریف نوں ویکھ کے جانچیا جا سکدا اے۔

ہُن فہرست اُتے غور کرو:
\[
f(N_1), f(N_2), f(N_3), \dots
\]
کیوں جے~$f$، !!{surjective} اے، ایس لئی~$\Bin^\omega$ دا ہر رکن
~$f$ دی کسے دلیل لئی قدر لازماً اے، تے سو فہرست وچ لازماً آوے گا۔
ایس لئی ایہ فہرست سارے~$\Bin^\omega$ دی گنتی وار فہرست اے۔

سو جے~$\Pow{\Nat}$، !!{enumerable} ہوندا، تاں~$\Bin^\omega$ وی
!!{enumerable} ہوندا۔ پر~$\Bin^\omega$، !!{nonenumerable} اے
(\olref[nen-alt]{thm:nonenum-bin-omega})۔ ایس لئی~$\Pow{\Nat}$،
!!{nonenumerable} اے۔
\end{proof}

%\begin{explain}
%تقلیل کس پاسے جاندی اے، ایس بارے بھلیکھا پینا سوکھا اے۔
%مثال لئی، کوئی !!a{surjective} فنکشن $g \colon \Bin^\omega \to X$
%ایہ \emph{نہیں} ثابت کردا کہ~$X$، !!{nonenumerable} اے۔ (فنکشن
%$g
%\colon \Bin^\omega \to \Bin$ اُتے غور کرو جیہڑا $g(s) = s(1)$
%نال تعریف اے، یعنی اوہ فنکشن جیہڑا $0$ تے $1$ دی اک لڑی نوں اوہدے
%پہلے !!{element} ول بھیجدا اے۔ ایہ ہر ہدف قدر لیندا اے، کیوں جے کجھ
%لڑیاں $0$ توں تے کجھ $1$ توں شروع ہوندیاں نیں۔ پر~$\Bin$ متناہی
%اے۔) ایہ وی یاد رکھو کہ فنکشن~$f$ دا ہر ہدف قدر لینا لازمی اے،
%نہیں تاں دلیل نہیں چلدی: $f(x_1)$، $f(x_2)$، \dots{} وچ~$Y$ دے
%سارے !!{element}s دے شامل ہون دی ضمانت نہیں ہو سکدی۔ مثال لئی،
%$h\colon \Nat \to
%\Bin^\omega$ جیہڑا ایویں تعریف اے:
%\[
%h(n) = \underbrace{000\dots0}_{\text{$n$ $0$'s}}
%\]
%اک فنکشن اے، پر~$\Nat$، !!{enumerable} اے۔
%\end{explain}

\begin{prob}\label{sfr:siz:red:prob:nat-nat}
  تقلیلی دلیل نال وکھاؤ کہ سیٹ~$X$، جیہڑا سارے فنکشناں
  $f\colon \Nat \to \Nat$ دا سیٹ اے، !!{nonenumerable} اے (اشارا:
  $X$ توں~$\Bin^\omega$ ول اک ہر ہدف قدر لین والا فنکشن دیو۔)
\end{prob}

\begin{prob}
تقلیلی دلیل نال وکھاؤ کہ قدرتی عدداں دیاں جوڑیاں دے سارے
\emph{سیٹاں دے} سیٹ، یعنی $\Pow{\Nat \times \Nat}$،
!!{nonenumerable} اے۔
\end{prob}

\begin{prob}
وکھاؤ کہ~$\Nat^\omega$، یعنی قدرتی عدداں دیاں لامتناہی لڑیاں دا سیٹ،
تقلیلی دلیل نال !!{nonenumerable} اے۔
\end{prob}

%\begin{prob}
%$P$، $\Nat$ توں سیٹ~$\{0\}$ ول فنکشناں دا سیٹ ہووے، تے~$Q$ مثبت
%صحیح عدداں دے سیٹ توں سیٹ~$\{0\}$ ول \emph{جزوی} فنکشناں دا سیٹ
%ہووے۔ وکھاؤ کہ~$P$، !!{enumerable} اے تے~$Q$ نہیں۔ (اشارا:
%$\Bin^\omega$ دی گنتی وار فہرست بناؤن دے مسئلے نوں~$Q$ دی گنتی وار
%فہرست بناؤن ول گھٹاؤ۔)
%\end{prob}

\begin{prob}
$S$، $\Nat$ توں سیٹ~$\{0,1\}$ ول سارے !!{surjection}s دا سیٹ
ہووے، یعنی~$S$ سارے !!{surjection}s
$f \colon \Nat \to \Bin$ اُتے مشتمل اے۔ وکھاؤ کہ~$S$،
!!{nonenumerable} اے۔
\end{prob}

\begin{prob}
وکھاؤ کہ سارے حقیقی عدداں دا سیٹ~$\Real$، !!{nonenumerable} اے۔
\end{prob}

\end{document}
